\documentclass[12pt,a4paper]{exam}
\usepackage[utf8]{inputenc}
\usepackage{amsmath,amssymb,amsfonts}
\usepackage{geometry}
\usepackage{enumitem}
\usepackage{graphicx}
\usepackage{hyperref}
\usepackage{xcolor}
\geometry{a4paper, top=2cm, bottom=2cm, left=2.5cm, right=2.5cm}
\hypersetup{colorlinks=true, linkcolor=blue, urlcolor=blue}
\renewcommand{\thequestion}{\textbf{\arabic{question}}}
\renewcommand{\questionlabel}{\thequestion.}
\pointname{ marks}
\pointsdroppedatright
\bracketedpoints
\setlength{\parindent}{0pt}
\setlength{\emergencystretch}{3em}
\allowdisplaybreaks
\newsavebox{\schmathbox}
\newcommand{\fitmath}[1]{%
  \sbox{\schmathbox}{$\displaystyle #1$}%
  \ifdim\wd\schmathbox>\linewidth
    \resizebox{\linewidth}{!}{\usebox{\schmathbox}}%
  \else
    \usebox{\schmathbox}%
  \fi}

\begin{document}

\thispagestyle{empty}
\begin{center}
  \textbf{\LARGE Diag} \\[0.3cm]
  \textbf{Time Allowed: 3 Hours} \hfill \textbf{Maximum Marks: 80} \\[0.3cm]
  \rule{\textwidth}{0.5pt}
\end{center}

\vspace{0.3cm}

\noindent\textbf{General Instructions:}
\begin{enumerate}
  \item {\normalfont\normalcolor This question paper contains 40 questions. All questions are compulsory.}
  \item {\normalfont\normalcolor Questions are grouped into sections A through E.}
  \item {\normalfont\normalcolor Use of calculators is NOT permitted.}
\end{enumerate}
\bigskip
\noindent\textbf{\large Section A: Multiple Choice \& Assertion-Reason (1 mark each)}

\begin{questions}
\question[1] A tangent $PQ$ at a point $P$ of a circle of radius $5 \text{ cm}$ meets a line through the centre $O$ at a point $Q$ so that $OQ = 13 \text{ cm}$. The length of $PQ$ is:
\begin{enumerate}[label=(\Alph*)]
  \item (A) 10 cm
  \item (B) 8 cm
  \item (C) 12 cm
  \item (D) 7 cm
\end{enumerate}
\vspace{0.15cm}

\question[1] If tangents $PA$ and $PB$ from a point $P$ to a circle with centre $O$ are inclined to each other at an angle of $80^\circ$, then $\angle POA$ is equal to:
\begin{enumerate}[label=(\Alph*)]
  \item $(A) 50^\circ$
  \item $(B) 60^\circ$
  \item $(C) 70^\circ$
  \item $(D) 80^\circ$
\end{enumerate}
\vspace{0.15cm}

\question[1] From a point $Q$, the length of the tangent to a circle is $24 \text{ cm}$ and the distance of $Q$ from the centre is $25 \text{ cm}$. The radius of the circle is:
\begin{enumerate}[label=(\Alph*)]
  \item (A) 9 cm
  \item (B) 12 cm
  \item (C) 15 cm
  \item (D) 7 cm
\end{enumerate}
\vspace{0.15cm}

\question[1] Two concentric circles are of radii $5 \text{ cm}$ and $3 \text{ cm}$. The length of the chord of the larger circle which touches the smaller circle is:
\begin{enumerate}[label=(\Alph*)]
  \item (A) 10 cm
  \item (B) 8 cm
  \item (C) 6 cm
  \item (D) 4 cm
\end{enumerate}
\vspace{0.15cm}

\question[1] The number of tangents that can be drawn from a point inside a circle is:
\begin{enumerate}[label=(\Alph*)]
  \item (A) 1
  \item (B) 2
  \item (C) 0
  \item (D) Infinite
\end{enumerate}
\vspace{0.15cm}

\question[1] If the circumference of a circle is $176\text{ cm}$, then its radius is:
\begin{enumerate}[label=(\Alph*)]
  \item (A) 14 cm
  \item (B) 21 cm
  \item (C) 28 cm
  \item (D) 35 cm
\end{enumerate}
\vspace{0.15cm}

\question[1] The area of a circle whose diameter is $14\text{ cm}$ is:
\begin{enumerate}[label=(\Alph*)]
  \item (A) 44 cm$^{2}$
  \item (B) 154 cm$^{2}$
  \item (C) 616 cm$^{2}$
  \item (D) 308 cm$^{2}$
\end{enumerate}
\vspace{0.15cm}

\question[1] The area of a quadrant of a circle with radius $7\text{ cm}$ is:
\begin{enumerate}[label=(\Alph*)]
  \item (A) 38.5 cm$^{2}$
  \item (B) 77 cm$^{2}$
  \item (C) 154 cm$^{2}$
  \item (D) 19.25 cm$^{2}$
\end{enumerate}
\vspace{0.15cm}

\question[1] If the perimeter of a semi-circular protractor is $36\text{ cm}$, then its diameter is:
\begin{enumerate}[label=(\Alph*)]
  \item (A) 7 cm
  \item (B) 21 cm
  \item (C) 28 cm
  \item (D) 14 cm
\end{enumerate}
\vspace{0.15cm}

\question[1] The area of a sector of a circle with radius $6\text{ cm}$ and central angle $60^\circ$ is:
\begin{enumerate}[label=(\Alph*)]
  \item $(A) 3\pi cm^{2}$
  \item $(B) 4\pi cm^{2}$
  \item $(C) 6\pi cm^{2}$
  \item $(D) 12\pi cm^{2}$
\end{enumerate}
\vspace{0.15cm}

\question[1] Assertion (A): The length of a tangent drawn from an external point to a circle is always equal to the radius of the circle.
Reason (R): The tangent at any point of a circle is perpendicular to the radius through the point of contact.
\begin{enumerate}[label=(\Alph*)]
  \item (A) Both Assertion (A) and Reason (R) are true and Reason (R) is the correct explanation of Assertion (A)
  \item (B) Both Assertion (A) and Reason (R) are true but Reason (R) is NOT the correct explanation of Assertion (A)
  \item (C) Assertion (A) is true but Reason (R) is false
  \item (D) Assertion (A) is false but Reason (R) is true
\end{enumerate}
\vspace{0.15cm}

\question[1] Assertion (A): If the distance between the centres of two circles of radii $r_1$ and $r_2$ is $d$, then they touch each other externally if $d = r_1 + r_2$.
Reason (R): The distance between the centres of two circles is the sum of their radii when they touch internally.
\begin{enumerate}[label=(\Alph*)]
  \item (A) Both Assertion (A) and Reason (R) are true and Reason (R) is the correct explanation of Assertion (A)
  \item (B) Both Assertion (A) and Reason (R) are true but Reason (R) is NOT the correct explanation of Assertion (A)
  \item (C) Assertion (A) is true but Reason (R) is false
  \item (D) Assertion (A) is false but Reason (R) is true
\end{enumerate}
\vspace{0.15cm}

\question[1] Assertion (A): The number of tangents that can be drawn from a point inside a circle is zero.
Reason (R): A tangent to a circle intersects the circle at only one point.
\begin{enumerate}[label=(\Alph*)]
  \item (A) Both Assertion (A) and Reason (R) are true and Reason (R) is the correct explanation of Assertion (A)
  \item (B) Both Assertion (A) and Reason (R) are true but Reason (R) is NOT the correct explanation of Assertion (A)
  \item (C) Assertion (A) is true but Reason (R) is false
  \item (D) Assertion (A) is false but Reason (R) is true
\end{enumerate}
\vspace{0.15cm}

\question[1] Assertion (A): The length of the tangent from a point $P$ at a distance of $10 \text{ cm}$ from the centre of a circle of radius $6 \text{ cm}$ is $8 \text{ cm}$.
Reason (R): In a right triangle, the square of the hypotenuse equals the sum of the squares of the other two sides.
\begin{enumerate}[label=(\Alph*)]
  \item (A) Both Assertion (A) and Reason (R) are true and Reason (R) is the correct explanation of Assertion (A)
  \item (B) Both Assertion (A) and Reason (R) are true but Reason (R) is NOT the correct explanation of Assertion (A)
  \item (C) Assertion (A) is true but Reason (R) is false
  \item (D) Assertion (A) is false but Reason (R) is true
\end{enumerate}
\vspace{0.15cm}

\question[1] Assertion (A): Two parallel tangents can be drawn to a circle at the ends of a diameter.
Reason (R): Tangents at the ends of a diameter are perpendicular to the diameter.
\begin{enumerate}[label=(\Alph*)]
  \item (A) Both Assertion (A) and Reason (R) are true and Reason (R) is the correct explanation of Assertion (A)
  \item (B) Both Assertion (A) and Reason (R) are true but Reason (R) is NOT the correct explanation of Assertion (A)
  \item (C) Assertion (A) is true but Reason (R) is false
  \item (D) Assertion (A) is false but Reason (R) is true
\end{enumerate}
\vspace{0.15cm}

\question[1] Assertion (A): The area of a circle with radius $7$ cm is $154$ cm$^2$. Reason (R): The area of a circle is given by $\pi r^2$, where $r$ is the radius.
\begin{enumerate}[label=(\Alph*)]
  \item (A) Both Assertion (A) and Reason (R) are true and Reason (R) is the correct explanation of Assertion (A)
  \item (B) Both Assertion (A) and Reason (R) are true but Reason (R) is NOT the correct explanation of Assertion (A)
  \item (C) Assertion (A) is true but Reason (R) is false
  \item (D) Assertion (A) is false but Reason (R) is true
\end{enumerate}
\vspace{0.15cm}

\question[1] Assertion (A): The circumference of a circle with diameter $14$ cm is $44$ cm. Reason (R): The circumference of a circle is given by $\pi d$, where $d$ is the diameter.
\begin{enumerate}[label=(\Alph*)]
  \item (A) Both Assertion (A) and Reason (R) are true and Reason (R) is the correct explanation of Assertion (A)
  \item (B) Both Assertion (A) and Reason (R) are true but Reason (R) is NOT the correct explanation of Assertion (A)
  \item (C) Assertion (A) is true but Reason (R) is false
  \item (D) Assertion (A) is false but Reason (R) is true
\end{enumerate}
\vspace{0.15cm}

\question[1] Assertion (A): The area of a semicircle with radius $r$ is $\frac{1}{2} \pi r^2$. Reason (R): The circumference of a semicircle is $\pi r$.
\begin{enumerate}[label=(\Alph*)]
  \item (A) Both Assertion (A) and Reason (R) are true and Reason (R) is the correct explanation of Assertion (A)
  \item (B) Both Assertion (A) and Reason (R) are true but Reason (R) is NOT the correct explanation of Assertion (A)
  \item (C) Assertion (A) is true but Reason (R) is false
  \item (D) Assertion (A) is false but Reason (R) is true
\end{enumerate}
\vspace{0.15cm}

\question[1] Assertion (A): If the perimeter and the area of a circle are numerically equal, then the radius of the circle is $2$ units. Reason (R): Circumference $2\pi r = \pi r^2$ implies $r=2$.
\begin{enumerate}[label=(\Alph*)]
  \item (A) Both Assertion (A) and Reason (R) are true and Reason (R) is the correct explanation of Assertion (A)
  \item (B) Both Assertion (A) and Reason (R) are true but Reason (R) is NOT the correct explanation of Assertion (A)
  \item (C) Assertion (A) is true but Reason (R) is false
  \item (D) Assertion (A) is false but Reason (R) is true
\end{enumerate}
\vspace{0.15cm}

\question[1] Assertion (A): The area of a sector of a circle with radius $r$ and central angle $\theta$ is $\frac{\theta}{360} \times \pi r^2$. Reason (R): The length of the arc of a sector is $\frac{\theta}{360} \times \pi r^2$.
\begin{enumerate}[label=(\Alph*)]
  \item (A) Both Assertion (A) and Reason (R) are true and Reason (R) is the correct explanation of Assertion (A)
  \item (B) Both Assertion (A) and Reason (R) are true but Reason (R) is NOT the correct explanation of Assertion (A)
  \item (C) Assertion (A) is true but Reason (R) is false
  \item (D) Assertion (A) is false but Reason (R) is true
\end{enumerate}
\vspace{0.15cm}

\end{questions}
\bigskip
\noindent\textbf{\large Section B: Very Short Answer (2 marks each)}

\begin{questions}
\question[2] Two concentric circles are of radii $5 \text{ cm}$ and $3 \text{ cm}$. Find the length of the chord of the larger circle which touches the smaller circle.
\vspace{0.15cm}

\question[2] If a tangent $PA$ drawn from an external point $P$ to a circle with center $O$ makes an angle of $30^\circ$ with the chord $AB$, where $AB$ is the diameter, find $\angle APB$.
\vspace{0.15cm}

\question[2] In the given figure, if $\angle OAB = 40^\circ$, find $\angle ACB$ where $O$ is the center of the circle.
\vspace{0.15cm}

\question[2] A quadrilateral $\text{ABCD}$ is drawn to circumscribe a circle. If $AB = 6 \text{ cm}$, $BC = 7 \text{ cm}$ and $CD = 4 \text{ cm}$, find the length of $AD$.
\vspace{0.15cm}

\question[2] Two tangents $TP$ and $TQ$ are drawn to a circle with center $O$ from an external point $T$. If $\angle PTQ = 70^\circ$, find $\angle POQ$.
\vspace{0.15cm}

\question[2] Find the area of a sector of a circle with radius $6 \text{ cm}$ if the angle of the sector is $60^\circ$. (Use $\pi = 3.14$)
\vspace{0.15cm}

\end{questions}
\bigskip
\noindent\textbf{\large Section C: Short Answer (3 marks each)}

\begin{questions}
\question[3] The angle of elevation of the top of a chimney from the top of a tower is $30^\circ$ and the angle of depression of the foot of the chimney from the top of the tower is $60^\circ$. If the tower is $40$ m high, find the height of the chimney.
\vspace{0.15cm}

\question[3] A quadrilateral $\text{ABCD}$ is drawn to circumscribe a circle. Prove that $AB + CD = AD + BC$.
\vspace{0.15cm}

\question[3] Two concentric circles are of radii $5\text{ cm}$ and $3\text{ cm}$. Find the length of the chord of the larger circle which touches the smaller circle.
\vspace{0.15cm}

\question[3] If a tangent $PA$ and $PB$ from a point $P$ to a circle with centre $O$ are inclined to each other at an angle of $80^\circ$, then find $\angle POA$.
\vspace{0.15cm}

\question[3] Prove that the lengths of tangents drawn from an external point to a circle are equal.
\vspace{0.15cm}

\question[3] A tangent $PQ$ at a point $P$ of a circle of radius $5\text{ cm}$ meets a line through the centre $O$ at a point $Q$ so that $OQ = 12\text{ cm}$. Find the length $PQ$.
\vspace{0.15cm}

\end{questions}
\bigskip
\noindent\textbf{\large Section D: Long Answer (4-5 marks each)}

\begin{questions}
\question[5] A straight highway leads to the foot of a tower. A man standing at the top of the tower observes a car at an angle of depression of $30^\circ$, which is approaching the foot of the tower with a uniform speed. Six seconds later, the angle of depression of the car is found to be $60^\circ$. Find the time taken by the car to reach the foot of the tower from this point.
\vspace{0.15cm}

\question[5] The angle of elevation of the top of a tower from a point on the ground is $30^\circ$. On walking $40$ metres towards the tower, the angle of elevation changes to $60^\circ$. Find the height of the tower and the distance of the initial point from the tower.
\vspace{0.15cm}

\question[5] A vertical tower stands on a horizontal plane and is surmounted by a vertical flagstaff of height $6$ metres. At a point on the plane, the angle of elevation of the bottom and top of the flagstaff are $30^\circ$ and $60^\circ$ respectively. Find the height of the tower.
\vspace{0.15cm}

\question[5] From the top of a $7$ m high building, the angle of elevation of the top of a cable tower is $60^\circ$ and the angle of depression of its foot is $45^\circ$. Determine the height of the tower.
\vspace{0.15cm}

\question[5] Prove that the lengths of tangents drawn from an external point to a circle are equal.
\vspace{0.15cm}

\question[5] $PQ$ is a chord of length $8\text{ cm}$ of a circle of radius $5\text{ cm}$. The tangents at $P$ and $Q$ intersect at a point $T$. Find the length of $TP$.
\vspace{0.15cm}

\question[5] A quadrilateral $\text{ABCD}$ is drawn to circumscribe a circle. Prove that $AB + CD = AD + BC$.
\vspace{0.15cm}

\question[5] Two concentric circles are of radii $5\text{ cm}$ and $3\text{ cm}$. Find the length of the chord of the larger circle which touches the smaller circle.
\vspace{0.15cm}

\end{questions}
\bigskip

\end{document}
